% Cover letter — Applied Intelligence submission.
% Build: latexmk -pdf cover_letter.tex
% Teks merah = keputusan/isian penulis sebelum dikirim.
\documentclass[11pt,a4paper]{article}
\usepackage[T1]{fontenc}
\usepackage[margin=2.5cm]{geometry}
\usepackage{xcolor}
\usepackage[hidelinks]{hyperref}
\usepackage{enumitem}
\setlength{\parindent}{0pt}
\setlength{\parskip}{6pt}
\newcommand{\todo}[1]{\textcolor{red!70!black}{[#1]}}

\begin{document}

\begin{flushright}
29 October 2026
\end{flushright}

The Editor-in-Chief\\
\emph{Applied Intelligence}\\
Springer Nature

Dear Editor-in-Chief,

On behalf of all authors, I am pleased to submit our original research
article entitled \textbf{``Sequential Hybrid Sampling with Informed
Undersampling for Imbalanced Drug--Target Interaction Prediction''} for
consideration for publication in \emph{Applied Intelligence}.

Drug--target interaction (DTI) prediction is a real-life problem in which
the positive class is extremely scarce, with positive rates between 6\%
and below 1\% on the standard benchmark. Resampling is the usual remedy,
yet most resampling studies in this field rely on a single train--test
split without statistical testing. Our manuscript evaluates a sequential
hybrid sampling framework, in which the minority class is oversampled
beyond the majority size and an undersampling stage then selects the
balancing subset by using ant colony optimization (ACO), simulated
annealing (SA), or random selection as an uninformed control. To our
knowledge, this is the first systematic evaluation of this
overproduce-and-choose resampling paradigm in DTI prediction.

We compared 13 methods on the four gold-standard Yamanishi datasets in
725 runs under a repeated multi-seed protocol with paired statistical
tests, a full factorial ablation of the ACO undersampler, and a
budget-matched SA comparator. The main findings are as follows.
\begin{itemize}[nosep]
  \item A single split is unreliable, in which the F1 score of one
        configuration ranged from 0.267 to 0.431 across splits.
  \item The benefit of resampling depends on the minority richness of the
        dataset, in which the best F1 score rises from 0.32 to 0.47 on the
        smallest dataset, while no method beats the baseline on the
        richest one.
  \item On that dataset, the ACO undersampler abstained in all 15 runs and
        returned the original training data, which we identify and verify
        as a built-in safety valve for optimization-based resampling.
  \item The budget-matched SA is statistically indistinguishable from ACO
        on every dataset ($p \geq 0.125$), so that the value lies in
        informed selection rather than in a particular metaheuristic,
        while informed selection costs up to about 460 times more than
        random selection.
\end{itemize}

We believe the manuscript fits the scope of \emph{Applied Intelligence}
as it applies metaheuristic search to a complex real-life problem and
reports its benefits, its failure modes, and its computational cost
with equal care. It also directly extends a line of work that the
journal has published, i.e., simulated-annealing-based undersampling
(SAUS), which serves as the basis of our SA comparator. All code, data
pipelines, and results are publicly available on GitHub and archived on
Zenodo (DOI \href{https://doi.org/10.5281/zenodo.23686130}{10.5281/zenodo.23686130}).

We confirm that this manuscript is original, has not been published
elsewhere, and is not under consideration by any other journal. All
authors have approved the manuscript and agree with its submission to
\emph{Applied Intelligence}. The authors declare no competing interests.

\todo{KEPUTUSAN PENULIS --- pengungkapan karya terkait. Panduan jurnal
meminta corresponding author menyebut materi terkait yang belum terbit.
Usulan kalimat: ``A related but distinct manuscript by our group, which
reports a single-split evaluation of an ACO-based sequential hybrid
sampling on one dataset, is currently under review at the IAES
International Journal of Artificial Intelligence. The present manuscript
differs in its datasets, protocol, comparators, and findings, and it
does not reproduce any text, table, or figure from that manuscript.''
Hapus blok ini bila diputuskan tidak relevan.}

\todo{KEPUTUSAN PENULIS --- pengungkapan penggunaan AI. Kebijakan
Springer Nature mewajibkan penggunaan LLM di luar copy editing
didokumentasikan di bagian Methods. Bila disetujui, tambahkan satu
kalimat di Methods dan kalimat berikut di sini: ``The use of a large
language model during the preparation of this work is described in the
Methods section, and the authors take full responsibility for the
content of the manuscript.''}

\textbf{Suggested reviewers.} \todo{Verifikasi independensi (bukan
kolaborator, bukan satu institusi), afiliasi terkini, dan email
institusional sebelum diisi. Kandidat dari literatur yang disitasi
naskah: penulis ESMOTE (Zhang et al., Neural Computing and
Applications, 2023), penulis SAUS (Chennuru and Timmappareddy, Applied
Intelligence, 2022), penulis FF-SMOTE (Kaur and Gosain, Applied
Artificial Intelligence, 2019), dan penulis ACOSampling (Yu et al.,
Neurocomputing, 2013). Panduan jurnal menganjurkan campuran negara dan
institusi.}

Thank you for considering our manuscript. We look forward to hearing
from you.

Sincerely,

\medskip
\textbf{Isman Kurniawan} (corresponding author)\\
School of Computing, Telkom University, Bandung, Indonesia\\
\href{mailto:ismankrn@telkomuniversity.ac.id}{ismankrn@telkomuniversity.ac.id}

\end{document}
